\chapter{Minimal structural information}\label{ch:minimal}

\Cref{thm:nerve} says that the dead-set pattern $D_A$ determines $\beta_2$. How much of $D_A$ must an observer see? This chapter measures, on the full six-label census, how much \emph{count-level} information about the nerve --- numbers of vertices, edges, triangles, Euler characteristics, numbers of hollow triangles --- leaves uncertain about whether $\beta_2>0$, and contrasts it with a \emph{structural} reduction that preserves the answer exactly: strong collapse of the nerve to its core, followed by one $H_1$ computation. The summary: counts are insufficient; the core is sufficient by a theorem; the core is not claimed to be minimal.

\section{Strong collapses and cores}
Let $\mathcal K$ be a finite simplicial complex. A vertex $v$ is \emph{dominated} by a vertex $w\neq v$ if every maximal simplex of $\mathcal K$ containing $v$ also contains $w$. Deleting a dominated vertex (together with all simplices containing it) is a \emph{strong collapse}; the result is a subcomplex that is a deformation retract of $\mathcal K$.
\begin{theorem}[strong collapses preserve homotopy type; \cite{BarmakMinian}]\label{thm:barmak}
If $v$ is dominated in $\mathcal K$ then $\mathcal K\setminus v\simeq\mathcal K$ (indeed $\mathcal K$ strong-collapses onto $\mathcal K\setminus v$). Iterating until no vertex is dominated yields a complex $\mathcal K_0$, the \emph{core} of $\mathcal K$, with $\mathcal K_0\simeq\mathcal K$; the isomorphism class of the core does not depend on the order of collapses.
\status{\Classical}{Barmak and Minian 2012}
\end{theorem}
\begin{corollary}[the core suffices]\label{cor:core-suffices}
For prefix-defined $A$ with $\emptyset\ne A\subsetneq S_6$, let $\mathcal C_A$ be the core of $N[D_A]$. Then $\beta_2(K_A)=\dim H_1(\mathcal C_A;\F)$ and $\beta_1(K_A)=\dim\Hh_2(\mathcal C_A;\F)$.
\status{\Analytic\ (\cref{thm:nerve} with \cref{thm:barmak}); re-verified by the manuscript on $160$ random regimes (all with $\beta_2=0$) and on all $290$ known positives, where $\beta_2$ from the definition of $K_A$ equals $\dim H_1$ of the nerve and of the core and is positive in every case (\texttt{ch12\_13\_nerve\_and\_core\_random}, \texttt{ch12\_13\_all\_290\_positives})}{companion paper, Section 6}
\end{corollary}
\begin{proof}
$N[D_A]\simeq Y_A$ by \cref{thm:nerve}, and the core is homotopy equivalent to the nerve by \cref{thm:barmak}. Homology is a homotopy invariant; the Betti formulas are \cref{thm:AD}.
\end{proof}

\section{The measurement}\label{sec:D21}
Let $\mathcal P_{21}$ be the set of distinct (rank-indexed) dead-set patterns arising in the exhaustive six-label census (labels in $\{1,\dots,16\}$): the census has $346{,}859$ regimes, of which $8{,}008$ are fully admissible (there $Y_A=\emptyset$ and $K_A$ is the whole boundary of the polytope) and are excluded, leaving $338{,}851$ regimes with $219{,}145$ distinct patterns, $151$ of them with $\beta_2>0$ (carried by the $158$ positive regimes). For each pattern the following observations were computed exactly: the numbers $V,E$ of vertices and edges of the nerve and the number of triangles; the Euler characteristic of its $2$-skeleton; the number of hollow triangles; the size of the strong-collapse core; and the answer to ``$H_1$ of the core''. For each observation the exact conditional entropy $H(Y\mid O)$ of $Y=\mathbf 1[\beta_2>0]$ was computed in the natural ensemble (uniform over patterns) and the balanced ensemble (positives and negatives with probability $1/2$ each).
\begin{table}[ht]\centering\small
\begin{tabular}{@{}lrrr@{}}\toprule
observation $O$ of the nerve & $H(Y|O)$ natural & $H(Y|O)$ balanced & classes\\ \midrule
$V=|D|$ & $0.0049$ & $0.5122$ & $57$\\
$(V,E)$ & $0.0037$ & $0.2418$ & $388$\\
Euler characteristic of the $2$-skeleton & $0.0020$ & $0.1073$ & $10{,}930$\\
number of hollow triangles & $0.0038$ & $0.1919$ & $2{,}639$\\
$(V,E,\#\text{triangles})$ & $0.0006$ & $0.0052$ & $20{,}595$\\
collapse-core size & $0.0009$ & $0.0319$ & $57$\\
(core size, hollow count) & $0.00002$ & $0.00007$ & $11{,}468$\\
collapse core, then $H_1$ & $0$ & $0$ & ---\\ \bottomrule
\end{tabular}
\caption{Exact residual uncertainty (bits) about $\beta_2>0$ over the six-label census, for count-level observations of the dead-set nerve (D21). The last row is a control: by \cref{cor:core-suffices} it computes $Y$ exactly. Values are inherited from the focused companion paper and were not recomputed.}\label{tab:D21}
\end{table}
\begin{observation}[no tested count-level statistic of the nerve is sufficient]\label{obs:collapse}
In \cref{tab:D21} every count-level row has positive residual entropy: counts of vertices, edges, triangles, hollow triangles, the Euler characteristic and even the core size, alone or in the pairs tested, leave positives and negatives in common classes. Only the structural reduction (collapse to the core followed by $H_1$) has zero residual, and by \cref{cor:core-suffices} it does so as a theorem and not as a measurement. We do not claim that the core is the smallest possible observation, nor any out-of-sample statement beyond the census; these entropies describe a finite ensemble.
Structural facts about the cores, in the census:
\begin{itemize}[nosep]
\item Positive patterns have small cores: of the $151$, the strong-collapse core sizes are $3$ (for $105$ patterns: the bare hollow triangle), $5$ ($23$), $6$ ($6$), $7$ ($1$), $8$ ($2$), $14$ ($5$), $19$ ($1$), $20$ ($3$), $21$ ($3$) and $33$ ($2$) vertices. \emph{The manuscript verifier recomputed this distribution exactly from the dead-set patterns of the $158$ census positives}, obtaining the same $151$ patterns and the same sizes (\texttt{ch13\_census\_positive\_patterns\_core\_sizes}).
\item $148{,}096$ of the negative patterns collapse to a single point; the cores of the remaining negatives reach $57$ vertices but have vanishing first homology (inherited).
\end{itemize}
\status{Table and negative-core statistics: \FiniteExact\ (inherited, \texttt{NOT\_RERUN}); positive-core distribution: \FiniteExact\ (re-derived)}{companion paper, Section 6; GIG-INFO v005}
\end{observation}
\begin{figure}[ht]\centering\includegraphics[width=\textwidth]{figures/fig8_collapse_core.pdf}
\caption{(a) The dead-set nerve of content $(2,3,4,5,8,15)$ at source $480$ ($11$ dead sets). (b) Its strong-collapse core: three vertices forming the boundary of a triangle (a hollow triangle), so $\beta_2=1$. Both drawn from the exact computation. (c) Core sizes of the $151$ positive dead-set patterns of the census; the histogram is hard-coded in the figure generator and has been \emph{reproduced independently} by the verifier (\cref{obs:collapse}).}\label{fig:collapse}
\end{figure}
\begin{remark}[what ``minimal'' means here]
The question ``what is the minimal structural information that determines whether a hole exists?'' has three different answers depending on the observer. (i) \emph{Without structure} (counts only), no tested statistic suffices (\cref{tab:D21}). (ii) \emph{With the dead-set pattern}, the answer is determined (\cref{thm:nerve}). (iii) \emph{With a homotopy-preserving reduction}, a hollow triangle in $105$ of the $151$ positive patterns \emph{is} the entire information. No universal optimality is claimed: a smaller sufficient observation of the pattern could exist, and no lower bound is proved.
\end{remark}

\section{The role of the hollow triangle}
In terms of ordinary combinatorics: a hollow triangle is a triple of forbidden prefix sets, pairwise compatible, with no prefix set related to all three. \Cref{obs:cycle-census} records that $288$ of $290$ positives contain a non-bounding one, and that hollow triangles are common among negatives ($707$ of $1000$) yet always bound there. The difference between a bounding and a non-bounding hollow triangle is a homological, not a local, condition: a dead set joined to all three vertices would bound the triangle, but triangles with no such cone point also bound through longer fillings ($6$ positive and $105$ negative cases). The first row of \cref{tab:D21} shows what the counts miss: the number of dead sets alone separates nothing, because the topology is carried by \emph{which} dead sets are present and how they are related by comparability, not by how many.
\begin{proposition}[counts cannot see the cycle --- a certified instance]\label{prop:counts-fail}
The two systems of \cref{prop:pair} at sources $2520$ and $3780$ have the same number ($5$) of dead sets and the same $N$, $E$, $Q$, $h_1$; their nerves have the same number of vertices ($5$), but differ in $\beta_2$ ($1$ vs $0$). The nerves have different edge sets, and the first one contains a non-bounding hollow triangle.
\status{\FiniteExact\ (recomputed, \texttt{figures/fig7\_data.json}, \texttt{ch13\_named\_systems\_cores})}{companion paper, Proposition 5}
\end{proposition}
